\documentclass[11pt]{article}

\usepackage[letterpaper,margin=1in]{geometry}
\usepackage[T1]{fontenc}
\usepackage[utf8]{inputenc}
\usepackage{lmodern}
\usepackage{microtype}
\usepackage{amsmath,amssymb,amsthm}
\usepackage{booktabs}
\usepackage{array}
\usepackage{xcolor}
\usepackage{enumitem}
\usepackage[hidelinks]{hyperref}

\definecolor{noteBlue}{HTML}{2E74B5}
\definecolor{noteGray}{HTML}{52606D}

\newtheorem{theorem}{Theorem}
\newtheorem{lemma}[theorem]{Lemma}
\newcommand{\fib}{F}
\newcommand{\R}{\mathbb{R}}

\title{\textbf{A Proof of the OEIS A105565 Discrepancy Bound}}
\author{
  Adam Haig\\[-2pt]
  \small human prompter and submitting author
  \and
  OpenAI Codex (AI system)\thanks{Codex is an AI system, not a legal person. Its
  inclusion in the byline credits its substantive proof and formalization work;
  Adam Haig is the human submitting author and is responsible for the deposit.}\\[-2pt]
  \small AI proof author and formalizer
}
\date{15 August 2026}

\begin{document}
\maketitle

\begin{abstract}
Let $a(n)$ indicate whether exactly five Fibonacci numbers have $n$ decimal
digits, and let $S(n)=\sum_{k=1}^{n}a(k)$.  OEIS A105565 conjectures a sharp
bounded discrepancy for $S(n)$.  We prove the stronger exact identity
\[
  S(n)=\lfloor n\alpha+\beta\rfloor-1 \qquad (n\geq 1),
\]
where $\alpha=\log(10)/\log(\varphi)-4$ and
$\beta=\log(5)/(2\log(\varphi))-1$.  The conjectured strict inequalities follow
immediately.  The proof accounts for the finite cutoff in the Formal
Conjectures definition and has been checked by Lean~4 against mathlib without
unproved assumptions.
\end{abstract}

\noindent\textbf{2020 Mathematics Subject Classification.}
11B39, 11B83, 68V15.

\section{The sequence and the conjecture}

Let $(\fib_k)_{k\geq 0}$ be the Fibonacci sequence with $\fib_0=0$ and
$\fib_1=1$.  For $n\geq 1$, let $a(n)=1$ when exactly five values $\fib_k$
have $n$ decimal digits, and let $a(n)=0$ otherwise.  Write
\[
 S(n)=\sum_{k=1}^{n}a(k), \qquad
 \varphi=\frac{1+\sqrt5}{2},
\]
and define
\[
 \alpha=\frac{\log 10}{\log\varphi}-4,
 \qquad
 \beta=\frac{\log 5}{2\log\varphi}-1.
\]
This is OEIS A105565, introduced by Reinhard Zumkeller.  Hans J.~H.~Tuenter
recorded the discrepancy conjecture in 2025 and also gave a fractional-part
description of $a(n)$ for $n>1$~\cite{oeis-a105565}.

\begin{theorem}\label{thm:main}
For every integer $n\geq 1$,
\[
 S(n)=\lfloor n\alpha+\beta\rfloor-1.
\]
Consequently,
\[
 \beta-2<S(n)-\alpha n<\beta-1.
\]
\end{theorem}

The result follows once the exact Fibonacci index cutoff is identified.  Two
boundary cases matter: no power $10^n$ is Fibonacci, and the logarithmic
cutoff is never an integer.  These facts also make the final inequalities
strict.

\section{The exact Fibonacci cutoff}

For $n\geq 1$, define
\[
 x_n=\log_{\varphi}(10^n\sqrt5).
\]
Binet's formula, with $\psi=(1-\sqrt5)/2$, gives
$\fib_k=(\varphi^k-\psi^k)/\sqrt5$.  Since $|\psi|<1$ and $\sqrt5>2$,
\begin{equation}\label{eq:error}
 \left|\fib_k-\frac{\varphi^k}{\sqrt5}\right|
   =\frac{|\psi|^k}{\sqrt5}<\frac12.
\end{equation}

\begin{lemma}[Nonintegral boundary]\label{lem:nonintegral}
For $n\geq 1$, $x_n$ is not an integer.
\end{lemma}

\begin{proof}
Suppose $x_n=k\in\mathbb Z$.  Positivity forces $k\geq 1$, and then
$\varphi^k=10^n\sqrt5$.  The identity
\[
 \varphi^k=\fib_k\varphi+\fib_{k-1}
\]
together with $\sqrt5=2\varphi-1$ yields
\begin{equation}\label{eq:rational-phi}
 (2\cdot 10^n-\fib_k)\varphi=10^n+\fib_{k-1}.
\end{equation}
The right side is positive, so the coefficient of $\varphi$ is nonzero.
Equation~\eqref{eq:rational-phi} would therefore make $\varphi$ rational, a
contradiction.
\end{proof}

\begin{lemma}[Powers of ten are not Fibonacci]\label{lem:not-power-ten}
For $n\geq 1$, there is no $k$ with $\fib_k=10^n$.
\end{lemma}

\begin{proof}
If equality held, then $2\mid\fib_k$ and $5\mid\fib_k$.  The identity
\[
 \gcd(\fib_r,\fib_s)=\fib_{\gcd(r,s)}
\]
with $r=3$ and $r=5$ implies $3\mid k$ and $5\mid k$, respectively.  Hence
$15\mid k$, so $\fib_{15}=610$ divides $\fib_k$.  In particular
$61\mid 10^n$.  Since $61$ is prime, this would imply $61\mid 10$, which is
impossible.
\end{proof}

\begin{lemma}[Cutoff equivalence]\label{lem:cutoff}
For $n\geq1$ and $k\geq0$,
\[
 \fib_k<10^n \quad\Longleftrightarrow\quad k<x_n.
\]
\end{lemma}

\begin{proof}
The inequality $k<x_n$ is equivalent to
$\varphi^k/\sqrt5<10^n$.  Equation~\eqref{eq:error} then gives
$\fib_k<10^n+1/2$.  Because $\fib_k$ and $10^n$ are integers,
$\fib_k\leq10^n$; equality is excluded by Lemma~\ref{lem:not-power-ten}.

Conversely, suppose $\fib_k<10^n$.  If $k\geq x_n$, then
Lemma~\ref{lem:nonintegral} gives $k>x_n$, so
$\varphi^k/\sqrt5>10^n$.  Equation~\eqref{eq:error} would imply
$\fib_k>10^n-1/2$, contradicting the integrality of $\fib_k$ and the strict
upper bound $\fib_k<10^n$.
\end{proof}

Set $c_n=\lfloor x_n\rfloor$.  Lemmas~\ref{lem:nonintegral} and
\ref{lem:cutoff} show that
\begin{equation}\label{eq:cutoff-floor}
 \fib_k<10^n \quad\Longleftrightarrow\quad k\leq c_n.
\end{equation}

\section{Reduction to a mechanical word}

Expanding the logarithm gives
\begin{equation}\label{eq:affine}
 x_n=4n+1+(n\alpha+\beta).
\end{equation}
Put $q_n=n\alpha+\beta$ and $b_n=\lfloor q_n\rfloor$.  Then
$c_n=4n+1+b_n$.

The elementary bounds $\varphi^4<10<\varphi^5$ give
\begin{equation}\label{eq:alpha-bounds}
 0<\alpha<1.
\end{equation}
Similarly, $\varphi^6<10\sqrt5<\varphi^7$, combined with
$x_1=5+\alpha+\beta$, gives
\begin{equation}\label{eq:alpha-beta-bounds}
 1<\alpha+\beta<2.
\end{equation}
In particular $q_n>0$.  Equations~\eqref{eq:alpha-bounds} and
\eqref{eq:alpha-beta-bounds} also give $q_n<n+2$, so $b_n<n+2$ and
\begin{equation}\label{eq:finite-cutoff}
 c_n<5n+10.
\end{equation}
This verifies that the finite search range $[0,5n+10)$ in the Formal
Conjectures definition contains every relevant Fibonacci index.  The formal
theorem therefore proves the repository definition itself, not merely an
unbounded reformulation.

For $n\geq2$, equation~\eqref{eq:cutoff-floor} shows that the indices of
Fibonacci numbers with exactly $n$ digits are precisely the integers in
$(c_{n-1},c_n]$.  Their number is
\begin{equation}\label{eq:count}
 c_n-c_{n-1}=4+(b_n-b_{n-1}).
\end{equation}
Because $0<\alpha<1$, the floor increment
$d_n=b_n-b_{n-1}$ is either $0$ or $1$.  Equation~\eqref{eq:count} shows that
exactly five $n$-digit Fibonacci numbers occur exactly when $d_n=1$.  Hence
\begin{equation}\label{eq:a-difference}
 a(n)=b_n-b_{n-1}\qquad(n\geq2).
\end{equation}

\begin{proof}[Proof of Theorem~\ref{thm:main}]
The initial value is $a(1)=0$.  By
equation~\eqref{eq:alpha-beta-bounds}, $b_1=\lfloor\alpha+\beta\rfloor=1$.
Summing equation~\eqref{eq:a-difference} therefore telescopes:
\[
 S(n)=b_n-1=\lfloor n\alpha+\beta\rfloor-1.
\]
Lemma~\ref{lem:nonintegral} and equation~\eqref{eq:affine} show that $q_n$ is
never an integer.  Consequently $q_n-1<b_n<q_n$.  Substituting
$q_n=n\alpha+\beta$ and $S(n)=b_n-1$ gives
\[
 \beta-2<S(n)-\alpha n<\beta-1,
\]
as required.
\end{proof}

\section{Machine-checked verification}

The argument was formalized in Lean~4 against the Google DeepMind Formal
Conjectures repository~\cite{formal-repo,formal-paper}.  The formalization
proves the exact bounded-filter sequence definition and the theorem
\texttt{OeisA105565.conjecture}.  It was checked against repository commit
\begin{center}
\small\texttt{2411d22e1bd550d050d0eac6c1fb379a76a3e7c5}
\end{center}
using Lean/mathlib 4.27.0.  The following checks completed successfully.

\begin{center}
\small
\begin{tabular}{>{\bfseries}p{0.16\linewidth}p{0.72\linewidth}}
\toprule
Check & Result \\
\midrule
File & \texttt{lake env lean FormalConjectures/OEIS/105565.lean} passed. \\
Module & \texttt{lake --wfail build FormalConjectures.OEIS.105565} completed
8,055 jobs successfully. \\
Axioms & \texttt{propext}, \texttt{Classical.choice}, and
\texttt{Quot.sound}; no \texttt{sorryAx}. \\
Sanity & No \texttt{sorry}, \texttt{admit}, \texttt{axiom}, or
\texttt{unsafe} declaration occurs in the target file. \\
\bottomrule
\end{tabular}
\end{center}

A kernel check does not replace editorial review of exposition or novelty, but
it certifies that the theorem follows from the displayed Lean definitions and
mathlib's accepted axioms~\cite{lean4,mathlib}.

\section*{Contribution and AI-use statement}

Adam Haig supplied the prompt, selected the conjecture and acceptance criteria,
authorized publication, and reviewed the resulting mathematical and formal
artifacts.  OpenAI Codex, an AI coding and reasoning system, performed the
majority of the mathematical work: it developed and refined the cutoff
argument, located relevant mathlib lemmas, generated and iteratively repaired
the Lean formalization, executed the verification checks reported above, and
drafted this note.  The byline credits both contributors while identifying
Codex explicitly as an AI system.  Haig is the human submitting author and
accepts responsibility for the deposit and for correcting any error that may
be found.

\section*{Acknowledgements}

The authors thank Reinhard Zumkeller for creating A105565, Hans J.~H.~Tuenter
for recording its formula and discrepancy conjecture, the OEIS community for
maintaining the entry, and the Formal Conjectures and mathlib contributors for
the machine-checkable framework.

\begin{thebibliography}{9}

\bibitem{oeis-a105565}
OEIS Foundation Inc.,
\emph{The On-Line Encyclopedia of Integer Sequences, A105565: Indicator for
five Fibonacci numbers with $n$ digits}.
\url{https://oeis.org/A105565} (accessed 15 August 2026).

\bibitem{spilker}
J.~Spilker,
\emph{Die Ziffern der Fibonacci-Zahlen},
Elemente der Mathematik \textbf{58} (2003), 26--33.
\url{https://doi.org/10.1007/PL00012537}.

\bibitem{formal-repo}
The Formal Conjectures Authors,
\emph{The Formal Conjectures Repository}, 2025--2026.
\url{https://github.com/google-deepmind/formal-conjectures}.

\bibitem{formal-paper}
M.~Firsching et al.,
\emph{Formal Conjectures: An Open and Evolving Benchmark for Verified
Discovery in Mathematics}, arXiv:2605.13171 (2026).
\url{https://arxiv.org/abs/2605.13171}.

\bibitem{lean4}
L.~de Moura and S.~Ullrich,
\emph{The Lean 4 theorem prover and programming language},
CADE 2021, LNCS 12699, 625--635.
\url{https://doi.org/10.1007/978-3-030-79876-5_37}.

\bibitem{mathlib}
The mathlib Community,
\emph{The Lean mathematical library},
Proceedings of CPP 2020, 367--381.
\url{https://doi.org/10.1145/3372885.3373824}.

\end{thebibliography}

\end{document}
